% answers_short.tex —— 2000 年·春季上海卷 · 简短答案速查
% 由 tools/make_short_answers.py 自 answers.tex 的 \gksolhead/\gkanswer 提取，
% 并与 2000/2000_数学真题_春季上海_答案要点.md 逐题比对后生成，勿手改（改源 MD 后重跑）。

\begin{gknotice}
\textbf{使用说明}：本节只给最终答案便于快速核对，完整推导见解析版。答案为编写者独立推导，\textbf{非官方标准答案}。
\end{gknotice}

\gkpart{填空题}
\begin{itemize}[leftmargin=2.6em]
  \item[\textbf{1.}] $-\dfrac{7}{25}$。
  \item[\textbf{2.}] $1$。
  \item[\textbf{3.}] $x=2$。
  \item[\textbf{4.}] $x=\dfrac{1}{a}$。
  \item[\textbf{5.}] $60^\circ$。
  \item[\textbf{6.}] $x+2y-1=0$。
  \item[\textbf{7.}] $\dfrac{3}{2}$。
  \item[\textbf{8.}] $45^\circ$。
  \item[\textbf{9.}] $5\sqrt{5}\,\mathrm{cm}$。
  \item[\textbf{10.}] $n=5$。
  \item[\textbf{11.}] $r=3$ 或 $r=7$。
  \item[\textbf{12.}] 例如 $\overline{Q}\cap P$。
\end{itemize}

\gkpart{选择题}
\begin{center}\begin{tabular}{ccccc}
\toprule
题号 & 13 & 14 & 15 & 16 \\
\midrule
答案 & B & D & A & A \\
\bottomrule
\end{tabular}\end{center}

\gkpart{解答题}
\begin{itemize}[leftmargin=2.6em]
  \item[\textbf{17.}] $f(x)=\begin{cases}x+2,&x\le-1,\\2-x^2,&-1<x<1,\\-x+2,&x\ge1.\end{cases}$；图见题面图稿。
  \item[\textbf{18.}] $z=\pm\left(\dfrac{\sqrt2}{2}+\dfrac{7\sqrt2}{2}i\right)$；$z=\dfrac{1+7i}{\sqrt2}$ 时 $m=0,2$，$z=-\dfrac{1+7i}{\sqrt2}$ 时 $m=0,-2$。
  \item[\textbf{19.}] 购买 $1$–$9$ 台时乙较省；购买 $10$ 台两者相同；购买 $11$ 台及以上时甲较省。
  \item[\textbf{20.}] (1) $a_n=6n-399$，$b_n=2\left(\dfrac9{10}\right)^{n-1}$；(2) $m=4,5,6,7,8$
  \item[\textbf{21.}] (1) $PA\perp$ 底面 $ABCD$，证明见解析；(2) 四棱锥体积为 $16$；(3) 混合积绝对值为 $48$，是四棱锥体积的 $3$ 倍；该混合积绝对值表示以三个向量为棱的平行六面体体积
  \item[\textbf{22.}] (1) $A=(1,3)$，$F=(1,1)$；当 $t\ne1$ 时直线 $TQ$ 为 $y=\dfrac{x-t}{1-t}$，当 $t=1$ 时为 $x=1$；(2) $t>\dfrac23$，$S=\dfrac{3t^2}{2(3t-2)}$，最小值为 $\dfrac43$；(3) 单调递增区间为 $\left(\dfrac43,+\infty\right)$
\end{itemize}
